\section{Applying Regularization Terms to the TGCE$_{SS}$ Loss
Function}\label{sec:regularizers}

The model produces two distinct outputs: segmentation masks
$\mathbf{\hat{Y}} = f(\mathbf{X};\theta)$ and reliability maps
$\{\Lambda_r(\mathbf{X};\theta)\}_{r=1}^R$. These outputs work in
tandem with the $TGCE_{SS}$ loss function described in Section
\ref{sec:proposed-loss-functions}. The loss function simultaneously guides the
learning of both the segmentation masks and reliability maps,
ensuring that the model learns to balance the contributions of
different annotators based on their estimated reliability.

The TGCE$_{SS}$ loss function is implemented differently for each reliability
approach, with specific adaptations to handle the different spatial dimensions
of the reliability maps. Naturally, the model itself container tunable
hyperparameters related to the \gls{GCE} - \gls{MAE} trade off behavior
(also known as $q$ parameter) and the unreliability tolerance $\tilde C$,
however, some additional hyperpamarameters were added with the purpose of
regularizing the overall loss for every experiment.

These additional regularization terms arise from the need to ensure
consistency in the internal representations of the model, since some
unregularized experiments showed that the model could sometimes converge to
low loss values but outputs in the reliability branch with no
consistency at all, since the model might assign extreme reliability
values to some labelers and not others. The added regularization terms help the
model converge to reliability values that are more consistent across
all labelers, that do not reach extreme values and that maintain high
reliability scores for scorers that did applied a label in the first
place and also were considered better labelers \footnote{For this,
  the Oxford PET dataset with emulated noisy labelers was very
  useful since the labeler reliability (related to noise levels in encoder
  layers) was very easy to control, making it possible
  to debug the model and note if reliability estimations levels were consistent
  with the labelers considered better. This was the main criteria to
design and select the regularization terms.}

\subsection{Centroid Penalty for Extreme Reliabilities}

Some optimization processes might favor a convergence to extreme reliability
estimations, by assigning either 0 or 1 to all labelers. To avoid this, a
regularization term is added to the loss function which penalizes the model
for assigning extreme reliability values and favoring those around a
centroid. High scaling values for this regularization term might
convergence to a very homogeneous reliability map across all
labelers, which is not desired either.

\begin{equation}
  \mathcal{L}_{centroid} = \alpha_{centroid} \cdot
  \mathbb{E}_{r \in R}[(\Lambda_r - 0.5)^2]
\end{equation}

The behavior of this regularization term is shown in Figure
\ref{fig:regularization_terms} (a).

\subsection{Entropy regularization}
In this case, entropy is useful to ensure that the reliability maps
are not too homogeneous. The entropy regularization term encourages
the model to assign more discriminative reliability values by penalizing
uniform distributions. This prevents the model from converging to
uninformative reliability estimates where all labelers appear equally
reliable, thus maintaining meaningful distinctions between more and
less trustworthy annotators.

\begin{equation}
  \mathcal{L}_{entropy} = \alpha_{entropy} \cdot \mathbb{E}_{r \in R}[\Lambda_r
    \log(1 + \Lambda_r)
  +(1 - \Lambda_r) \log(1 + (1 - \Lambda_r))]
\end{equation}

The behavior of this regularization term is shown in Figure
\ref{fig:regularization_terms} (b).

\subsection{Missing labels penalty}
In situations were the datasets have a low labeling rate, the model might
take an optimization shortcut by assigning high reliability values to
labelers which did not label at all, since the element-wise product with the
mask would be close to zero, which favors a minimization target.
To avoid this, a term was added to the loss function which penalizes
the model for assigning high reliability values to unlabeled patches,
as follows:

\begin{equation}
  \mathcal{L}_{\lambda} = \alpha_{\lambda} \cdot \max \left(0,
    \sum_{\forall r \in R}
    \left(\Lambda_r - \mu_r\right) \cdot
  e^{\frac{\Lambda_r - \mu_r - 1}{\beta}}\right)
\end{equation}

Where $\Lambda_r$ is the estimated reliability of the labeler $r$ from the
reliability branch for a certain patch, $\mu_r$ is the encoded
presence of the label mask for labeler $r$ for the same patch,
$\beta$ is a hyperparameter which control the penalty strength for
high and low estimated reliabilities, respectively.

The behavior of this regularization term is shown in Figure
\ref{fig:regularization_terms} (c).

\begin{figure}[H]
  \centering
  \begin{subfigure}[b]{0.48\textwidth}
    \begin{tikzpicture}
      \begin{axis}[
          axis lines=middle,
          xlabel={$\Lambda_r$},
          ylabel={$\mathcal{L}_{centroid}^r(\Lambda_r)$},
          width=\textwidth, height=6cm,
          domain=0:1,
          samples=200,
          smooth,
          thick,
          ymin=0, ymax=0.25,
          xtick={0,0.5,1.0},
          ytick={0,0.1,0.2,0.3},
          every axis x label/.style={at={(ticklabel* cs:1.05)}, anchor=west},
          every axis y label/.style={at={(ticklabel* cs:1.05)}, anchor=south},
        ]
        % Plot: (x - 0.5)^2
        \addplot[red,thick] {(x - 0.5)^2};
      \end{axis}
    \end{tikzpicture}
    \caption{$\mathcal{L}_{centroid}^r(\Lambda_r) = (\Lambda_r - 0.5)^2$}
    \label{fig:centroid_penalty}
  \end{subfigure}
  \hfill
  \begin{subfigure}[b]{0.48\textwidth}
    \begin{tikzpicture}
      \begin{axis}[
          axis lines=middle,
          xlabel={$\Lambda_r$},
          ylabel={$\mathcal{L}_{entropy}^r(\Lambda_r)$},
          width=\textwidth, height=6cm,
          domain=0:1,
          samples=200,
          smooth,
          thick,
          ymin=0, ymax=1.4,
          xtick={0,0.5,1.0},
          ytick={0,0.5,1.0,1.5},
          every axis x label/.style={at={(ticklabel* cs:1.05)}, anchor=west},
          every axis y label/.style={at={(ticklabel* cs:1.05)}, anchor=south},
        ]
        % Plot: x*ln(1+x) + (1-x)*ln(1+(1-x))
        \addplot[green,thick] {x*ln(1+x) + (1-x)*ln(1+(1-x))};
      \end{axis}
    \end{tikzpicture}
    \caption{$\mathcal{L}_{entropy}(\Lambda_r) = \Lambda_r \log(1 +
    \Lambda_r) + (1 - \Lambda_r) \log(1 + (1 - \Lambda_r))$}
    \label{fig:entropy_regularization}
  \end{subfigure}

  \vspace{0.5cm}

  \begin{subfigure}[b]{0.48\textwidth}
    \centering
    \begin{tikzpicture}
      \begin{axis}[
          axis lines=middle,
          xlabel={$z=\Lambda_r - \mu_r$},
          ylabel={$\mathcal{L}_{\lambda}^r(z)$},
          width=\textwidth, height=6cm,
          domain=-2:4,
          samples=400,
          smooth,
          thick,
          ymin=0, ymax=5,
          xtick={-2,-1,0,1,2,3,4},
          ytick={0,1,2,3,4,5},
          every axis x label/.style={at={(ticklabel* cs:1.05)}, anchor=west},
          every axis y label/.style={at={(ticklabel* cs:1.05)}, anchor=south},
        ]
        % Parameters
        \def\alphaval{0.2}
        \def\betaval{1.0}
        % Plot: max(0, expression)
        \addplot[blue,thick] {max(0, (x)*exp((x -
        1)/\betaval))};
      \end{axis}
    \end{tikzpicture}
    \caption{$\mathcal{L}_{\lambda}^r(z) = \max(0,
    z \cdot e^{\frac{z-1}{\beta}})$}
    \label{fig:missing_labels_penalizing}
  \end{subfigure}

  \caption{Visualization of the three regularization terms used in the
    TGCE loss function:
    (a) centroid penalty for extreme reliabilities, (b) entropy regularization,
    and (c) missing labels penalty.
    Each subplot shows
    how the respective regularization term behaves as a function of its
  input variables.}
  \label{fig:regularization_terms}
\end{figure}

This function applies no penalty for cases in which  $\Lambda_r <
\mu_r$, meaning no penalty for assigning any reliability value
(ranging from 0 to 1) to a labeler who actually labeled a given patch.
On the other hand, a penalty is indeed applied in cases where
$\Lambda_r > \mu_r$, meaning a penalty for assigning high reliability
values to a labeler which did not label at all. The proportion on how
high and how quickly (in the space of difference $\Lambda_r - \mu_r$)
the penalty is applied is controlled by the hyperparameters $\alpha_{\lambda}$
and $\beta$. The behavior of this term is shown in Figure
\ref{fig:regularization_terms} (d).

These regularization terms lead to the overall loss function being in the form:

\begin{align}
  L_{TGCEreg} &= L_{TGCEss} \\
  &+\alpha_{centroid} \cdot \mathbb{E}_r[(\Lambda_r - 0.5)^2] \nonumber \\
  &+\alpha_{entropy} \cdot \mathbb{E}_r[\Lambda_r \log(1 + \Lambda_r)
  +(1 - \Lambda_r) \log(1 + (1 - \Lambda_r))] \nonumber \\
  &+\alpha_{\lambda} \cdot \max \left(0, \sum_{\forall r \in R}
    \left(\Lambda_r - \mu_r\right) \cdot
  e^{\frac{\Lambda_r - \mu_r - 1}{\beta}}\right) \nonumber
\end{align}

For the scalar reliability approach, the loss function operates on
global reliability values per annotator. The reliability maps are
expanded to match the spatial dimensions of the predictions through
tiling operations. This approach is computationally efficient but
provides the same reliability value across all spatial locations.

The feature-based reliability approach requires additional processing
to handle the spatial dimensions of the bottleneck features. The
reliability maps are resized using bilinear interpolation to match
the spatial dimensions of the predictions. This approach provides
coarse-grained spatial reliability information while maintaining
computational efficiency.

The pixel-wise reliability approach operates directly on full-resolution
reliability maps, requiring no spatial transformations. This approach
provides the most detailed spatial reliability information but requires
more computational resources and is more prone to overfitting in
sparse datasets (low labeling rate).
